\anexo{Diseño de Entrevistas en Profundidad}

Como complemento cualitativo de la encuesta, se diseñaron tres entrevistas semiestructuradas a usuarios reales de Mercado Libre: dos vendedores y un comprador, seleccionados por muestreo intencional entre personas con actividad verificable en la plataforma (no se buscan perfiles profesionales ni expertos, sino usuarios comunes, en línea con el público objetivo de la herramienta). El propósito es profundizar en el porqué de los patrones cuantitativos observados en la encuesta: la detección tardía de tendencias (88,9\% combinado), la alta disposición a usar una herramienta automatizada (68,9\%) y el peso de las señales sociales en la decisión de compra (47,5\%). Cada entrevista dura entre 20 y 30 minutos, se registra con consentimiento informado y se analiza mediante síntesis temática. A continuación se presenta el diseño completo de las tres guías; los hallazgos se integrarán a esta sección una vez realizadas las entrevistas.

\subsection*{Entrevista 1 -- Vendedor de Mercado Libre (foco: la problemática)}

\emph{Objetivo}: entender el proceso real de decisión de stock y dimensionar el costo de detectar tendencias tarde.

\begin{enumerate}
	\item ¿Hace cuánto vendés en Mercado Libre, en qué rubro y qué tan seguido incorporás productos nuevos?
	\item Contame la última vez que sumaste un producto nuevo: ¿cómo decidiste que convenía?
	\item ¿Qué fuentes usás hoy para saber qué se está vendiendo o qué está creciendo?
	\item ¿Cuánto tiempo por semana le dedicás a investigar tendencias o productos para stockear?
	\item ¿Te pasó de llegar tarde a una tendencia? Contame el caso y qué consecuencia tuvo.
	\item Si esa información la hubieras tenido dos semanas antes, ¿qué habrías hecho distinto?
	\item ¿Cómo distinguís una moda pasajera de una tendencia que se va a sostener?
	\item ¿Qué información te falta hoy y no encontrás en ningún lado?
	\item ¿Probaste alguna herramienta o método para esto (planillas, apps, seguir cuentas)? ¿Qué te aportó?
	\item ¿Qué tendría que tener una herramienta para que la uses todas las semanas?
\end{enumerate}

\subsection*{Entrevista 2 -- Vendedor de Mercado Libre (foco: validación de la solución)}

\emph{Objetivo}: validar la comprensión y utilidad del dashboard y del score de tendencia mostrando el sistema en funcionamiento.

\begin{enumerate}
	\item ¿Qué vendés, hace cuánto, y cómo decidís hoy qué stockear?
	\item Antes de ver la herramienta: si una aplicación te dijera ``este producto está creciendo'', ¿qué necesitarías saber para creerle?
	\item (Mostrando el dashboard) ¿Qué entendés que muestra esta pantalla? (sin explicación previa)
	\item El score va de 0 a 100: ¿qué esperarías de un producto con 80 frente a uno con 40?
	\item ¿Qué harías mañana con este ranking si fuera tuyo?
	\item De los componentes del score -- aparece seguido, se mantiene en el tiempo, posición en el catálogo, precio estable -- ¿cuál te importa más? ¿Agregarías alguno?
	\item ¿Preferís entrar a mirar el panel o recibir una alerta cuando algo empieza a crecer? ¿Por qué medio?
	\item ¿Qué le falta a esta herramienta para que la uses en serio en tu operación?
	\item ¿La usarías gratis? ¿Pagarías por ella? ¿Qué valor mensual te parecería razonable?
	\item ¿Qué es lo que menos te convence de lo que viste?
\end{enumerate}

\subsection*{Entrevista 3 -- Comprador de Mercado Libre}

\emph{Objetivo}: profundizar el hallazgo de la encuesta sobre señales sociales (47,5\%) y la percepción de repetición de resultados destacados (59,1\%).

\begin{enumerate}
	\item ¿Qué tan seguido comprás en Mercado Libre y qué tipo de productos?
	\item Contame tu última compra de algo nuevo o novedoso: ¿cómo llegaste a ese producto?
	\item Cuando algo te lo recomiendan o lo ves en redes (TikTok/Instagram), ¿cuánto influye en que lo termines comprando?
	\item ¿Confiás más en una recomendación de conocidos, en las reseñas, o en lo que la plataforma te muestra primero?
	\item Cuando buscás algo, ¿sentís que te aparecen siempre los mismos productos destacados? ¿Eso te ayuda o te molesta?
	\item ¿Alguna vez compraste algo ``que estaba de moda'' y te arrepentiste? ¿Qué señales no viste a tiempo?
	\item ¿Comparás precios entre vendedores o vas al primero que aparece? ¿Cómo lo hacés?
	\item ¿Qué te daría más confianza para descubrir productos nuevos (historial de precios, ver si viene creciendo, opiniones)?
	\item Si pudieras ver que un producto ``viene subiendo hace dos semanas'' o que su precio está estable, ¿cambiaría tu decisión?
	\item ¿Qué es lo que más te frustra hoy al comprar algo nuevo en la plataforma?
\end{enumerate}
